\documentclass[11pt,a4paper]{article}
\usepackage[utf8]{inputenc}
\usepackage[spanish,es-tabla]{babel}
\usepackage{amsmath,amsfonts,amssymb}
\usepackage{geometry}
\usepackage{graphicx}
\usepackage{xcolor}

\usepackage{enumitem}
\usepackage{tikz}
\usetikzlibrary{positioning,shapes,arrows.meta,calc}

\geometry{margin=2.5cm}

\title{\textbf{Soluci\'on Examen Final 2021\\TAF ISC\\IMT Atlantique}}
\author{}
\date{}

\begin{document}

\maketitle

\section*{Ejercicio 1 - Estimaci\'on de canal / Ecualizaci\'on}

\textbf{Contexto del problema:}

Sistema de transmisi\'on con las siguientes caracter\'isticas:
\begin{itemize}
    \item S\'imbolos de informaci\'on: $a_k \in \{+1, -1, +i, -i\}$ (QPSK)
    \item Ruido: $b_n$ blanco gaussiano, centrado, varianza $\sigma^2$ conocida
    \item Canal discreto equivalente: $H(z) = \sum_{k=0}^{L-1} h_k z^{-k}$
\end{itemize}

\textbf{Se\~nal recibida:}
\[
y_n = \sum_{k=0}^{L-1} h_k a_{n-k} + b_n
\]

\subsection*{Pregunta 1: Intercorrelaci\'on se\~nal recibida con secuencia de entrenamiento}

\textbf{Enunciado:} Suponiendo los s\'imbolos $a_k$ conocidos (secuencia de entrenamiento), calcular la intercorrelaci\'on del se\~nal recibida $y_k$ con la secuencia de s\'imbolos $a_k$:
\[
R_{ya}(u) = \sum_k y(k) a^*(k-u)
\]

\vspace{0.3cm}
\textbf{Soluci\'on paso a paso:}

\textbf{Paso 1: Sustituir la expresi\'on de $y_n$}

\textit{\textbf{¿Qu\'e estamos haciendo?}} Queremos correlacionar la se\~nal recibida con la secuencia que enviamos (que conocemos). Esta t\'ecnica se llama \textbf{matched filtering} y es fundamental en comunicaciones.

\textit{Motivaci\'on:} Si el canal es $h$ y enviamos $a$, recibimos $y = h * a + ruido$. Al correlacionar $y$ con $a$, estamos "extrayendo" informaci\'on sobre $h$.

Partimos de la definici\'on de intercorrelaci\'on:
\[
R_{ya}(u) = \sum_k y_k a^*(k-u)
\]

\textit{Nota:} $a^*$ es el conjugado complejo (importante para QPSK que es complejo).

Sustituimos la expresi\'on de la se\~nal recibida $y_k = \sum_{l=0}^{L-1} h_l a_{k-l} + b_k$:
\[
R_{ya}(u) = \sum_k \left(\sum_{l=0}^{L-1} h_l a_{k-l} + b_k\right) a^*(k-u)
\]

\textbf{Paso 2: Separar en dos t\'erminos}

\[
R_{ya}(u) = \sum_k \sum_{l=0}^{L-1} h_l a_{k-l} a^*(k-u) + \sum_k b_k a^*(k-u)
\]

\textbf{Paso 3: Analizar el t\'ermino de ruido}

\textit{\textbf{¿Por qu\'e el ruido desaparece?}} Esta es la magia de la correlaci\'on:

Como $b_k$ es ruido blanco (independiente) y $a_k$ es nuestra secuencia conocida:
\[
\mathbb{E}\left[\sum_k b_k a^*(k-u)\right] = \sum_k \mathbb{E}[b_k] \cdot a^*(k-u) = 0
\]

porque $\mathbb{E}[b_k] = 0$ (ruido centrado).

\textit{En la pr\'actica:} Con $N$ muestras grandes, por la ley de los grandes n\'umeros:
\[
\frac{1}{N}\sum_k b_k a^*(k-u) \approx 0
\]

\textit{Analog\'{\i}a:} Es como promediar mediciones ruidosas - el ruido aleatorio se cancela, pero la se\~nal sistem\'atica (el canal) permanece.

\textbf{Paso 4: Cambiar orden de sumatorias}

Para el t\'ermino de se\~nal:
\[
R_{ya}(u) \approx \sum_{l=0}^{L-1} h_l \sum_k a_{k-l} a^*(k-u)
\]

\textbf{Paso 5: Reconocer la autocorrelaci\'on}

\textit{\textbf{Concepto clave - Autocorrelaci\'on:}} Mide cu\'anto se parece una secuencia a s\'{\i} misma desplazada.

La suma interna es exactamente la autocorrelaci\'on de nuestra secuencia de entrenamiento:
\[
\sum_k a_{k-l} a^*(k-u) = R_{aa}(u-l)
\]

\textit{Interpretaci\'on:}
\begin{itemize}
    \item Si $u = l$: $R_{aa}(0)$ = energ\'{\i}a de la secuencia (m\'aximo)
    \item Si $u \neq l$: mide similitud entre versiones desplazadas
    \item Ideal: $R_{aa}(u) \approx \delta(u)$ (cero excepto en $u=0$)
\end{itemize}

\textbf{Resultado final:}

\[
\boxed{R_{ya}(u) = \sum_{l=0}^{L-1} h_l R_{aa}(u-l)}
\]

Esta es la **convoluci\'on** entre la respuesta del canal $\{h_l\}$ y la autocorrelaci\'on de la secuencia de entrenamiento $R_{aa}$:

\[
\boxed{R_{ya} = h * R_{aa}}
\]

\textbf{Interpretaci\'on:}

\begin{itemize}
    \item La intercorrelaci\'on $R_{ya}$ contiene informaci\'on sobre el canal $h$
    \item Esta informaci\'on est\'a ''distorsionada'' por la autocorrelaci\'on de la secuencia de entrenamiento
    \item Para recuperar $h$, necesitamos que $R_{aa}$ sea lo m\'as parecida a una delta de Dirac
\end{itemize}



\subsection*{Pregunta 2: M\'etodo de estimaci\'on de la respuesta impulsional}

\textbf{Enunciado:} Suponiendo $R_{aa}(u) = \delta(u)$ con $\delta(u)$ la funci\'on de Dirac, proponer un m\'etodo de estimaci\'on de la respuesta impulsional.

\vspace{0.3cm}
\textbf{Soluci\'on paso a paso:}

\textbf{Paso 1: Aplicar la condici\'on ideal}

Si $R_{aa}(u) = \delta(u)$, entonces de la pregunta anterior:
\[
R_{ya}(u) = \sum_{l=0}^{L-1} h_l R_{aa}(u-l) = \sum_{l=0}^{L-1} h_l \delta(u-l)
\]

\textbf{Paso 2: Simplificar usando la propiedad de la delta}

\textit{\textbf{La propiedad m\'agica de la delta de Dirac:}} Funciona como un "selector" - extrae solo el valor en un punto espec\'{\i}fico.

La delta de Dirac tiene la propiedad fundamental:
\[
\sum_l f(l) \delta(u-l) = f(u)
\]

\textit{Explicaci\'on:} $\delta(u-l)$ es cero para todo $l \neq u$, y vale 1 en $l = u$. Por tanto, la suma solo "recoge" el valor $f(u)$.

Aplicando esto a nuestro caso:
\[
R_{ya}(u) = \sum_{l=0}^{L-1} h_l \delta(u-l) = h_u
\]

\textit{\textbf{¡Resultado extraordinario!}} La intercorrelaci\'on nos da \textbf{directamente} los coeficientes del canal.

\textbf{Paso 3: M\'etodo de estimaci\'on}

\[
\boxed{\hat{h}_u = R_{ya}(u) = \sum_k y_k a^*(k-u)}
\]

\textbf{Algoritmo de estimaci\'on por correlaci\'on cruzada:}

\begin{enumerate}
    \item \textbf{Transmitir:} Secuencia de entrenamiento $\{a_0, a_1, \ldots, a_{N-1}\}$ con $R_{aa}(u) \approx \delta(u)$
    
    \item \textbf{Recibir:} Secuencia $\{y_0, y_1, \ldots, y_{N-1}\}$
    
    \item \textbf{Calcular:} Para cada $u = 0, 1, \ldots, L-1$:
    \[
    \hat{h}_u = \frac{1}{N} \sum_{k=0}^{N-1} y_k a^*(k-u)
    \]
    
    \item \textbf{Obtener:} Estimaci\'on del canal $\{\hat{h}_0, \hat{h}_1, \ldots, \hat{h}_{L-1}\}$
\end{enumerate}

\textbf{Paso 4: Implementaci\'on pr\'actica}

En forma matricial:
\[
\begin{bmatrix}
y_0 \\
y_1 \\
y_2 \\
\vdots \\
y_{N-1}
\end{bmatrix}
=
\begin{bmatrix}
a_0 & 0 & 0 & \cdots & 0 \\
a_1 & a_0 & 0 & \cdots & 0 \\
a_2 & a_1 & a_0 & \cdots & 0 \\
\vdots & \vdots & \vdots & \ddots & \vdots \\
a_{N-1} & a_{N-2} & a_{N-3} & \cdots & a_{N-L}
\end{bmatrix}
\begin{bmatrix}
h_0 \\
h_1 \\
h_2 \\
\vdots \\
h_{L-1}
\end{bmatrix}
+
\begin{bmatrix}
b_0 \\
b_1 \\
b_2 \\
\vdots \\
b_{N-1}
\end{bmatrix}
\]

O en forma compacta: $\mathbf{y} = \mathbf{A} \mathbf{h} + \mathbf{b}$

La estimaci\'on por correlaci\'on es:
\[
\hat{\mathbf{h}} = (\mathbf{A}^H \mathbf{A})^{-1} \mathbf{A}^H \mathbf{y}
\]

Si $R_{aa} = \delta$, entonces $\mathbf{A}^H \mathbf{A} = N \mathbf{I}$, simplificando a:
\[
\hat{\mathbf{h}} = \frac{1}{N} \mathbf{A}^H \mathbf{y}
\]

\textbf{Ventajas del m\'etodo:}

\begin{itemize}
    \item \textbf{Simple:} Solo requiere multiplicaciones y sumas
    \item \textbf{Robusto:} Promedia el ruido sobre $N$ muestras
    \item \textbf{No iterativo:} Soluci\'on directa, no requiere convergencia
    \item \textbf{\'Optimo:} Si $R_{aa} = \delta$, es el estimador de m\'inimos cuadrados
\end{itemize}

\textbf{Precisi\'on de la estimaci\'on:}

La varianza del error de estimaci\'on es:
\[
\text{Var}(\hat{h}_u - h_u) = \frac{\sigma^2}{N}
\]

Mayor $N$ (secuencia m\'as larga) $\to$ mejor estimaci\'on.



\subsection*{Pregunta 3: Elecci\'on de secuencia de entrenamiento}

\textbf{Enunciado:} Tiene dos opciones de secuencias de entrenamiento. La primera tiene funci\'on de autocorrelaci\'on dada por Figura 1 (pico central grande, l\'obulos laterales peque\~nos). La segunda tiene autocorrelaci\'on dada por Figura 2 (pico central m\'as bajo, l\'obulos laterales m\'as grandes). ¿Cu\'al elegir y por qu\'e?

\vspace{0.3cm}
\textbf{Soluci\'on paso a paso:}

\textbf{Paso 1: Recordar el objetivo ideal}

\textit{\textbf{¿Qu\'e hace una buena secuencia de entrenamiento?}} Debe ser "ortogonal a s\'{\i} misma" cuando se desplaza.

Queremos que la autocorrelaci\'on $R_{aa}(u)$ sea lo m\'as parecida posible a un impulso $\delta(u)$:
\[
R_{aa}(u) = \begin{cases}
\text{M\'aximo (toda la energ\'{\i}a)} & u = 0 \\
0 \text{ (sin similitud)} & u \neq 0
\end{cases}
\]

\textit{\textbf{¿Por qu\'e?}} Si recordamos la ecuaci\'on:
\[
R_{ya}(u) = \sum_{l=0}^{L-1} h_l R_{aa}(u-l)
\]

Si $R_{aa}$ es un impulso perfecto, entonces $R_{ya}(u) = h_u$ directamente. Si no, los coeficientes $h_l$ se "mezclan" entre s\'{\i}.

\textbf{Paso 2: Analizar la Secuencia 1}

\textbf{Caracter\'isticas (seg\'un descripci\'on t\'ipica):}
\begin{itemize}
    \item Pico central: Alto y estrecho
    \item L\'obulos laterales: Peque\~nos (cercanos a cero)
    \item Ejemplo: Secuencia de Barker, c\'odigo de Gold
\end{itemize}

\textbf{Autocorrelaci\'on ideal:}
\[
R_{aa,1}(u) \approx \begin{cases}
A & u = 0 \\
\epsilon & u \neq 0, \quad \epsilon \ll A
\end{cases}
\]

\textbf{Paso 3: Analizar la Secuencia 2}

\textbf{Caracter\'isticas:}
\begin{itemize}
    \item Pico central: M\'as bajo
    \item L\'obulos laterales: M\'as grandes
    \item Ejemplo: Secuencia aleatoria corta
\end{itemize}

\textbf{Autocorrelaci\'on:}
\[
R_{aa,2}(u) = \begin{cases}
B & u = 0, \quad B < A \\
\delta & u \neq 0, \quad \delta > \epsilon
\end{cases}
\]

\textbf{Paso 4: Impacto en la estimaci\'on del canal}

Recordemos: $R_{ya}(u) = \sum_{l=0}^{L-1} h_l R_{aa}(u-l)$

\textbf{Con Secuencia 1 (l\'obulos peque\~nos):}
\[
\hat{h}_u = R_{ya}(u) \approx A \cdot h_u + \epsilon \sum_{l \neq u} h_l
\]

El t\'ermino de interferencia es peque\~no: $\epsilon \sum_{l \neq u} h_l \ll A \cdot h_u$

\textbf{Con Secuencia 2 (l\'obulos grandes):}
\[
\hat{h}_u = R_{ya}(u) \approx B \cdot h_u + \delta \sum_{l \neq u} h_l
\]

El t\'ermino de interferencia es significativo: $\delta \sum_{l \neq u} h_l$ no es despreciable.

\textbf{Paso 5: Consecuencias}

\textbf{Secuencia 1:}
\begin{itemize}
    \item $\checkmark$ Estimaci\'on precisa de cada coeficiente $h_u$
    \item $\checkmark$ M\'inima interferencia entre taps del canal
    \item $\checkmark$ Mejor relaci\'on se\~nal-interferencia
    \item Desventaja: M\'as dif\'icil de generar
\end{itemize}

\textbf{Secuencia 2:}
\begin{itemize}
    \item $\\times$ Cada $\hat{h}_u$ contiene contribuciones de otros $h_l$
    \item $\\times$ Sesgo en la estimaci\'on
    \item $\\times$ Error cuadr\'atico medio mayor
    \item Ventaja: M\'as f\'acil de generar
\end{itemize}

\textbf{Decisi\'on:}

\[
\boxed{\text{Elegir Secuencia 1 (l\'obulos laterales peque\~nos)}}
\]

\textbf{Justificaci\'on cuantitativa:}

Error cuadr\'atico medio de estimaci\'on:
\[
\text{MSE} = \mathbb{E}[|\hat{h}_u - h_u|^2]
\]

\textbf{Secuencia 1:}
\[
\text{MSE}_1 \approx \frac{\sigma^2}{A^2 N} + \frac{\epsilon^2}{A^2} \sum_{l \neq u} |h_l|^2
\]

\textbf{Secuencia 2:}
\[
\text{MSE}_2 \approx \frac{\sigma^2}{B^2 N} + \frac{\delta^2}{B^2} \sum_{l \neq u} |h_l|^2
\]

Como $\epsilon \ll \delta$, tenemos $\text{MSE}_1 \ll \text{MSE}_2$.

\textbf{Ejemplos de secuencias ideales:}

\begin{enumerate}
    \item \textbf{C\'odigo de Barker de longitud 13:}
    \[
    \{1,1,1,1,1,-1,-1,1,1,-1,1,-1,1\}
    \]
    L\'obulos laterales: m\'aximo $\pm 1/13 = 0.077$
    
    \item \textbf{Secuencias de Zadoff-Chu:}
    Autocorrelaci\'on perfecta (l\'obulos = 0 para desplazamientos c\'iclicos)
    
    \item \textbf{C\'odigos PN (Pseudo-Noise):}
    Para longitud $N$: l\'obulos $\approx 1/\sqrt{N}$
\end{enumerate}

\textbf{Conclusi\'on:}

Una buena secuencia de entrenamiento debe tener autocorrelaci\'on lo m\'as parecida a un impulso, con pico central alto y l\'obulos laterales m\'inimos. Esto minimiza la interferencia entre los coeficientes estimados del canal.



\subsection*{Pregunta 4: Ecualizador Zero-Forcing - Ventajas e inconvenientes}

\textbf{Enunciado:} Suponiendo conocer el canal discreto equivalente, se implementa el ecualizador:
\[
G(z) = \frac{1}{H(z)}
\]
Citar las ventajas e inconvenientes de este ecualizador.

\vspace{0.3cm}
\textbf{Soluci\'on:}

\textbf{Ventajas del ecualizador Zero-Forcing (ZF):}

\textbf{1. Elimina completamente la ISI (Ventaja principal)}

\textit{\textbf{Idea central del ZF:}} "Forzar a cero" la ISI invirtiendo exactamente el canal.

En el dominio de la frecuencia (transformada Z):
\[
Y(z) = H(z) X(z) + N(z)
\]

\textit{Se\~nal deseada + ruido}

Despu\'es de aplicar el ecualizador $G(z) = 1/H(z)$:
\[
\hat{X}(z) = G(z) Y(z) = \frac{1}{H(z)} [H(z) X(z) + N(z)] = X(z) + \frac{N(z)}{H(z)}
\]

\textit{\textbf{Observaci\'on crucial:}} El t\'ermino $H(z) X(z)$ se cancela perfectamente ($\frac{H(z)}{H(z)} = 1$), pero el ruido se transforma en $N(z)/H(z)$.

La se\~nal recuperada en el tiempo:
\[
\hat{x}_n = x_n + \text{ruido ecualizado}
\]

$\checkmark$ **ISI = 0** (en ausencia de ruido) - recuperaci\'on perfecta de la se\~nal

\textit{Analog\'{\i}a:} Si grabas audio en una habitaci\'on con eco y conoces exactamente el eco, puedes "revertirlo" para recuperar el audio original.

\textbf{2. Implementaci\'on conceptualmente simple}

\begin{itemize}
    \item Si $H(z) = \sum_{k=0}^{L-1} h_k z^{-k}$
    \item Entonces $G(z) = 1/H(z)$ es un filtro IIR
    \item O aproximar con FIR de longitud suficiente
\end{itemize}

\textbf{3. \'Optimo en ausencia de ruido}

Para canales sin ruido, ZF recupera perfectamente la se\~nal transmitida.

\vspace{0.5cm}
\textbf{Inconvenientes del ecualizador Zero-Forcing:}

\textbf{1. Amplificaci\'on catastr\'ofica del ruido (INCONVENIENTE PRINCIPAL)}

\textit{\textbf{El problema fundamental del ZF:}} Al invertir el canal, amplificamos enormemente el ruido en las frecuencias donde el canal es d\'ebil.

\textbf{An\'alisis en frecuencia:}

Recordemos que el ruido despu\'es del ecualizador es:
\[
N_{out}(f) = \frac{N(f)}{H(f)}
\]

\textit{\textbf{¿Qu\'e significa esto?}} Si el canal aten\'ua mucho una frecuencia (digamos $H(f_0) = 0.01$), el ecualizador debe amplificar por $1/0.01 = 100$ para compensar. ¡Pero esto amplifica el ruido tambi\'en!

Densidad espectral de potencia del ruido de salida:
\[
S_{N,out}(f) = S_{N,in}(f) \cdot \left|\frac{1}{H(f)}\right|^2 = \frac{N_0}{|H(f)|^2}
\]

\textbf{\textcolor{red}{Problema cr\'itico:}} Si $|H(f)| \to 0$ (desvanecimiento profundo o "notch"), entonces $|N_{out}(f)|^2 \to \infty$

\textbf{Ejemplo num\'erico detallado:}

\textit{\textbf{Escenario 1:}} Canal con atenuaci\'on moderada
\[
H(f_0) = 0.5 \quad \text{(atenuaci\'on del 50\% = -6 dB)}
\]

Amplificaci\'on del ruido:
\[
\frac{1}{|H(f_0)|^2} = \frac{1}{0.25} = 4 \quad \text{(+6 dB de ruido)}
\]

A\'un manejable.

\textit{\textbf{Escenario 2:}} Canal con desvanecimiento profundo
\[
H(f_1) = 0.1 \quad \text{(atenuaci\'on del 90\% = -20 dB)}
\]

Amplificaci\'on del ruido:
\[
\frac{1}{|H(f_1)|^2} = \frac{1}{0.01} = 100 \quad \text{(+20 dB de ruido)}
\]

\textcolor{red}{¡El ruido se amplifica 100 veces!}

\textit{\textbf{Escenario 3:}} Notch severo (casi nulo)
\[
H(f_2) = 0.01 \quad \text{(atenuaci\'on del 99\% = -40 dB)}
\]

Amplificaci\'on del ruido:
\[
\frac{1}{|H(f_2)|^2} = \frac{1}{0.0001} = 10,000 \quad \text{(+40 dB)}
\]

\textcolor{red}{\textbf{¡Catastr\'ofico! El ruido domina completamente la se\~nal.}}

\textbf{SNR despu\'es del ecualizador:}
\[
\text{SNR}_{out} = \frac{E_s}{\mathbb{E}[|N_{out}|^2]} = \frac{E_s}{N_0 \int \frac{1}{|H(f)|^2} df}
\]

Si el canal tiene notches (ceros espectrales), la integral diverge $\to$ SNR $\to$ 0.

\textbf{2. Inestabilidad num\'erica}

Si $H(z)$ tiene ceros cerca o fuera del c\'irculo unitario, $G(z) = 1/H(z)$ puede ser:
\begin{itemize}
    \item Inestable (polos fuera del c\'irculo unitario)
    \item No causal (polos en $z = 0$)
    \item Requiere aproximaciones que introducen errores
\end{itemize}

\textbf{3. No es \'optimo en sentido MMSE}

El ecualizador ZF minimiza la ISI pero ignora el ruido. El ecualizador \'optimo debe balancear:
\begin{itemize}
    \item Reducci\'on de ISI
    \item Amplificaci\'on de ruido
\end{itemize}

\textbf{4. Requiere conocimiento exacto del canal}

Errores en la estimaci\'on de $H(z)$ se traducen en:
\begin{itemize}
    \item ISI residual
    \item Mayor amplificaci\'on de ruido
    \item Posible inestabilidad
\end{itemize}

\textbf{5. Sensible a canales no-m\'inima fase}

Para canales no-m\'inima fase (ceros fuera del c\'irculo unitario), la inversa puede no ser realizable causalmente.

\vspace{0.5cm}
\textbf{Comparaci\'on cuantitativa:}

\begin{center}
\begin{tabular}{|l|c|c|}
\hline
\textbf{Caracter\'istica} & \textbf{ZF} & \textbf{MMSE} \\
\hline
ISI & 0 & Peque\~na \\
Amplificaci\'on ruido & Alta & Moderada \\
SNR salida & Baja & Alta \\
Complejidad & Media & Media \\
Robustez & Baja & Alta \\
\hline
\end{tabular}
\end{center}

\textbf{Cu\'ando usar ZF:}

\begin{itemize}
    \item Canal con respuesta plana (sin desvanecimientos profundos)
    \item SNR muy alto (ruido despreciable)
    \item Aplicaciones donde ISI = 0 es cr\'itico
\end{itemize}

\textbf{Alternativas preferibles:}

\begin{enumerate}
    \item \textbf{Ecualizador MMSE:}
    \[
    G_{MMSE}(z) = \frac{H^*(z)}{|H(z)|^2 + \sigma_n^2/\sigma_s^2}
    \]
    Balancea ISI y ruido
    
    \item \textbf{Ecualizador DFE:}
    Cancela ISI post-cursor sin amplificar ruido
    
    \item \textbf{Ecualizador ML (Viterbi):}
    \'Optimo pero computacionalmente costoso
\end{enumerate}

\textbf{Conclusi\'on:}

El ecualizador ZF es conceptualmente simple y elimina toda la ISI, pero su amplificaci\'on catastr\'ofica del ruido lo hace inadecuado para la mayor\'ia de aplicaciones pr\'acticas, especialmente en canales con desvanecimientos selectivos en frecuencia.



\newpage

\subsection*{Pregunta 5: Trellis del ecualizador ML para N=3}

\textbf{Enunciado:} Se quiere utilizar un ecualizador ML. En el caso donde $y_n = h_0 a_n + h_1 a_{n-1} + b_n$, con $h_0$ y $h_1$ conocidos, dibujar el trellis para un bloque de tama\~no $N = 3$. Se supone que el trellis comienza en el estado 1 (puede elegir este estado).

\vspace{0.3cm}
\textbf{Soluci\'on paso a paso:}

\textbf{Paso 1: Definir los estados del sistema}

\textit{\textbf{¿Qu\'e es un "estado" en este contexto?}} Es la informaci\'on del pasado que necesitamos para predecir la salida actual.

\textit{\textbf{¿Qu\'e informaci\'on necesitamos?}} Miremos la ecuaci\'on:
\[
y_n = h_0 a_n + h_1 a_{n-1} + b_n
\]

Para predecir $y_n$ dado $a_n$, necesitamos conocer $a_{n-1}$. Por tanto:
\[
\text{Estado en el instante } n: \quad S_n = a_{n-1}
\]

\textit{Interpretaci\'on:} El estado "recuerda" el s\'imbolo anterior (la memoria del canal).

Como $a_{n-1} \in \{+1, -1, +i, -i\}$ (QPSK tiene 4 s\'imbolos), hay exactamente 4 estados posibles:
\begin{itemize}
    \item Estado 1: $a_{n-1} = +1$ (parte real positiva del plano complejo)
    \item Estado 2: $a_{n-1} = -1$ (parte real negativa)
    \item Estado 3: $a_{n-1} = +i$ (parte imaginaria positiva)
    \item Estado 4: $a_{n-1} = -i$ (parte imaginaria negativa)
\end{itemize}

\textit{\textbf{Regla general:}} N\'umero de estados = $M^L$ donde $M$ = tama\~no del alfabeto, $L$ = memoria del canal.

\textbf{Paso 2: Definir las transiciones}

Desde cada estado, hay 4 transiciones posibles (seg\'un el valor de $a_n \in \{+1, -1, +i, -i\}$).

\textbf{Valor esperado en cada transici\'on:}

Para transici\'on de estado $S_j$ (donde $a_{n-1} = s_j$) al estado $S_k$ (donde $a_n = s_k$):
\[
y_{esperado} = h_0 s_k + h_1 s_j
\]

\textbf{Paso 3: Tabla de transiciones}

\begin{center}
\small
\begin{tabular}{|c|c|c|c|}
\hline
\textbf{Estado origen} & \textbf{$a_n$} & \textbf{Estado destino} & \textbf{$y_{esperado}$} \\
\hline
$S_1$ ($a_{n-1}=+1$) & $+1$ & $S_1$ & $h_0(+1) + h_1(+1)$ \\
$S_1$ ($a_{n-1}=+1$) & $-1$ & $S_2$ & $h_0(-1) + h_1(+1)$ \\
$S_1$ ($a_{n-1}=+1$) & $+i$ & $S_3$ & $h_0(+i) + h_1(+1)$ \\
$S_1$ ($a_{n-1}=+1$) & $-i$ & $S_4$ & $h_0(-i) + h_1(+1)$ \\
\hline
$S_2$ ($a_{n-1}=-1$) & $+1$ & $S_1$ & $h_0(+1) + h_1(-1)$ \\
$S_2$ ($a_{n-1}=-1$) & $-1$ & $S_2$ & $h_0(-1) + h_1(-1)$ \\
$S_2$ ($a_{n-1}=-1$) & $+i$ & $S_3$ & $h_0(+i) + h_1(-1)$ \\
$S_2$ ($a_{n-1}=-1$) & $-i$ & $S_4$ & $h_0(-i) + h_1(-1)$ \\
\hline
$S_3$ ($a_{n-1}=+i$) & $+1$ & $S_1$ & $h_0(+1) + h_1(+i)$ \\
$S_3$ ($a_{n-1}=+i$) & $-1$ & $S_2$ & $h_0(-1) + h_1(+i)$ \\
$S_3$ ($a_{n-1}=+i$) & $+i$ & $S_3$ & $h_0(+i) + h_1(+i)$ \\
$S_3$ ($a_{n-1}=+i$) & $-i$ & $S_4$ & $h_0(-i) + h_1(+i)$ \\
\hline
$S_4$ ($a_{n-1}=-i$) & $+1$ & $S_1$ & $h_0(+1) + h_1(-i)$ \\
$S_4$ ($a_{n-1}=-i$) & $-1$ & $S_2$ & $h_0(-1) + h_1(-i)$ \\
$S_4$ ($a_{n-1}=-i$) & $+i$ & $S_3$ & $h_0(+i) + h_1(-i)$ \\
$S_4$ ($a_{n-1}=-i$) & $-i$ & $S_4$ & $h_0(-i) + h_1(-i)$ \\
\hline
\end{tabular}
\end{center}

\textbf{Paso 4: Dibujar el trellis para N=3}

\begin{center}
\begin{tikzpicture}[scale=1.1, every node/.style={scale=0.75}]
% Título
\node[above] at (5.5, 6.5) {\textbf{Trellis del Ecualizador ML (QPSK, N=3)}};

% Ejes temporales
\node[below] at (0, -0.5) {$n=0$};
\node[below] at (3.5, -0.5) {$n=1$};
\node[below] at (7, -0.5) {$n=2$};
\node[below] at (10.5, -0.5) {$n=3$};

% Estados en cada instante
\foreach \x in {0, 3.5, 7, 10.5} {
    \node[draw, circle, fill=blue!20, minimum size=0.7cm] (s1\x) at (\x, 5) {$S_1$};
    \node[draw, circle, fill=red!20, minimum size=0.7cm] (s2\x) at (\x, 3.3) {$S_2$};
    \node[draw, circle, fill=green!20, minimum size=0.7cm] (s3\x) at (\x, 1.6) {$S_3$};
    \node[draw, circle, fill=yellow!20, minimum size=0.7cm] (s4\x) at (\x, 0) {$S_4$};
}

% Leyenda de estados
\node[left=0.3cm, align=right] at (s10.west) {\tiny $a_{n-1}=+1$};
\node[left=0.3cm, align=right] at (s20.west) {\tiny $a_{n-1}=-1$};
\node[left=0.3cm, align=right] at (s30.west) {\tiny $a_{n-1}=+i$};
\node[left=0.3cm, align=right] at (s40.west) {\tiny $a_{n-1}=-i$};

% Inicio en S1
\node[draw, star, star points=5, fill=orange, minimum size=0.5cm] at (-0.4, 5) {};
\node[left=0.1cm] at (-0.6, 5) {\tiny Inicio};

% Transiciones n=0 a n=1 (desde S1)
\draw[-latex, thick, blue] (s10) -- (s13.5) node[midway, above, sloped] {\tiny $a_1=+1$};
\draw[-latex, thick, red] (s10) -- (s23.5) node[midway, above, sloped] {\tiny $a_1=-1$};
\draw[-latex, thick, green!70!black] (s10) -- (s33.5) node[midway, above, sloped] {\tiny $a_1=+i$};
\draw[-latex, thick, orange] (s10) -- (s43.5) node[midway, above, sloped] {\tiny $a_1=-i$};

% Transiciones n=1 a n=2 (desde todos los estados)
% Desde S1
\draw[-latex, thin, blue] (s13.5) -- (s17);
\draw[-latex, thin, red] (s13.5) -- (s27);
\draw[-latex, thin, green!70!black] (s13.5) -- (s37);
\draw[-latex, thin, orange] (s13.5) -- (s47);

% Desde S2
\draw[-latex, thin, blue] (s23.5) -- (s17);
\draw[-latex, thin, red] (s23.5) -- (s27);
\draw[-latex, thin, green!70!black] (s23.5) -- (s37);
\draw[-latex, thin, orange] (s23.5) -- (s47);

% Desde S3
\draw[-latex, thin, blue] (s33.5) -- (s17);
\draw[-latex, thin, red] (s33.5) -- (s27);
\draw[-latex, thin, green!70!black] (s33.5) -- (s37);
\draw[-latex, thin, orange] (s33.5) -- (s47);

% Desde S4
\draw[-latex, thin, blue] (s43.5) -- (s17);
\draw[-latex, thin, red] (s43.5) -- (s27);
\draw[-latex, thin, green!70!black] (s43.5) -- (s37);
\draw[-latex, thin, orange] (s43.5) -- (s47);

% Transiciones n=2 a n=3 (similar)
% Desde S1
\draw[-latex, thin, blue] (s17) -- (s110.5);
\draw[-latex, thin, red] (s17) -- (s210.5);
\draw[-latex, thin, green!70!black] (s17) -- (s310.5);
\draw[-latex, thin, orange] (s17) -- (s410.5);

% Desde S2
\draw[-latex, thin, blue] (s27) -- (s110.5);
\draw[-latex, thin, red] (s27) -- (s210.5);
\draw[-latex, thin, green!70!black] (s27) -- (s310.5);
\draw[-latex, thin, orange] (s27) -- (s410.5);

% Desde S3
\draw[-latex, thin, blue] (s37) -- (s110.5);
\draw[-latex, thin, red] (s37) -- (s210.5);
\draw[-latex, thin, green!70!black] (s37) -- (s310.5);
\draw[-latex, thin, orange] (s37) -- (s410.5);

% Desde S4
\draw[-latex, thin, blue] (s47) -- (s110.5);
\draw[-latex, thin, red] (s47) -- (s210.5);
\draw[-latex, thin, green!70!black] (s47) -- (s310.5);
\draw[-latex, thin, orange] (s47) -- (s410.5);

\end{tikzpicture}
\end{center}

\textbf{Nota:} Cada transici\'on tiene 4 ramas saliendo (una por cada s\'imbolo QPSK posible).

\textbf{Paso 5: Algoritmo de Viterbi - Encontrando el camino \'optimo}

\textit{\textbf{El problema:}} Hay $4^3 = 64$ secuencias posibles de 3 s\'imbolos QPSK. ¿C\'omo encontrar la mejor sin probarlas todas?

\textit{\textbf{Soluci\'on de Viterbi:}} Usa \textbf{programaci\'on din\'amica} - descompone el problema grande en subproblemas m\'as peque\~nos.

\textbf{Concepto 1: M\'etrica de rama (branch metric)}

Para cada transici\'on, medimos qu\'e tan bien "encaja" con lo que realmente recibimos:
\[
\lambda_n = |y_n - y_{esperado}|^2 = |y_n - h_0 a_n - h_1 a_{n-1}|^2
\]

\textit{Interpretaci\'on:} Es la distancia cuadr\'atica entre:
\begin{itemize}
    \item Lo que recibimos: $y_n$
    \item Lo que predice esa transici\'on: $h_0 a_n + h_1 a_{n-1}$
\end{itemize}

Cuanto menor sea $\lambda_n$, mejor es esa hip\'otesis.

\textbf{Procedimiento:}

\begin{enumerate}
    \item \textbf{Inicializaci\'on ($n=0$):}
    \begin{itemize}
        \item Empezamos en estado $S_1$: $\Gamma_0(S_1) = 0$
        \item Otros estados: $\Gamma_0(S_i) = \infty$ para $i \neq 1$
    \end{itemize}
    
    \item \textbf{Para cada instante $n = 1, 2, 3$:}
    
    Para cada estado destino $S_j$:
    \begin{itemize}
        \item Calcular m\'etricas de las 4 ramas que llegan desde cada estado origen $S_i$:
        \[
        \Gamma_n^{(i \to j)} = \Gamma_{n-1}(S_i) + |y_n - h_0 a_n^{(j)} - h_1 a_{n-1}^{(i)}|^2
        \]
        
        \item Seleccionar la rama con m\'etrica m\'inima (survivor):
        \[
        \Gamma_n(S_j) = \min_i \Gamma_n^{(i \to j)}
        \]
        
        \item Guardar: decisi\'on del s\'imbolo $a_n$ y puntero al estado anterior
    \end{itemize}
    
    \item \textbf{Decisi\'on final ($n=3$):}
    \begin{itemize}
        \item Seleccionar estado con menor m\'etrica acumulada:
        \[
        S^* = \arg\min_{j} \Gamma_3(S_j)
        \]
        
        \item Trazar hacia atr\'as (backtracking) el camino \'optimo
        
        \item Extraer secuencia: $\{\hat{a}_1, \hat{a}_2, \hat{a}_3\}$
    \end{itemize}
\end{enumerate}

\textbf{Complejidad:}

\begin{itemize}
    \item Estados: $M = 4$ (QPSK)
    \item Ramas por estado: $M = 4$
    \item Operaciones por etapa: $M \times M = 16$ c\'alculos de m\'etrica + $M = 4$ comparaciones
    \item Total para $N=3$: $3 \times 16 = 48$ m\'etricas + $3 \times 4 = 12$ comparaciones
\end{itemize}

\textbf{Ventajas sobre DFE:}

\begin{itemize}
    \item \'Optimo en sentido ML
    \item No hay propagaci\'on de errores
    \item Decisiones globales sobre la secuencia completa
\end{itemize}

\textbf{Desventaja:}

Complejidad crece exponencialmente con la longitud de memoria del canal $L$.



\subsection*{Pregunta 6: N\'umero de estados para canal de memoria L=3}

\textbf{Enunciado:} Si ahora las muestras recibidas se expresan:
\[
y_n = h_0 a_n + h_1 a_{n-1} + h_2 a_{n-2} + b_n
\]
Dar el n\'umero de estados del nuevo trellis correspondiente.

\vspace{0.3cm}
\textbf{Soluci\'on paso a paso:}

\textbf{Paso 1: Identificar la memoria del canal}

La muestra $y_n$ depende de:
\begin{itemize}
    \item $a_n$ (s\'imbolo actual)
    \item $a_{n-1}$ (s\'imbolo anterior)
    \item $a_{n-2}$ (s\'imbolo 2 instantes atr\'as)
\end{itemize}

\textbf{Longitud de memoria:} $L = 2$ (depende de 2 s\'imbolos pasados)

\textbf{Paso 2: Definir el estado}

El estado debe contener toda la informaci\'on necesaria para predecir $y_n$:
\[
S_n = (a_{n-1}, a_{n-2})
\]

Es decir, el estado es un **par** de s\'imbolos anteriores.

\textbf{Paso 3: Contar los estados posibles}

Para QPSK, cada s\'imbolo puede tomar $M = 4$ valores: $\{+1, -1, +i, -i\}$

El estado es un par $(a_{n-1}, a_{n-2})$ donde:
\begin{itemize}
    \item $a_{n-1} \in \{+1, -1, +i, -i\}$: 4 opciones
    \item $a_{n-2} \in \{+1, -1, +i, -i\}$: 4 opciones
\end{itemize}

\textbf{N\'umero total de estados:}
\[
\boxed{\text{N\'umero de estados} = M^L = 4^2 = 16}
\]

\textbf{Paso 4: Listar todos los estados}

\begin{center}
\begin{tabular}{|c|c|c|}
\hline
\textbf{Estado} & \textbf{$a_{n-1}$} & \textbf{$a_{n-2}$} \\
\hline
$S_1$ & $+1$ & $+1$ \\
$S_2$ & $+1$ & $-1$ \\
$S_3$ & $+1$ & $+i$ \\
$S_4$ & $+1$ & $-i$ \\
\hline
$S_5$ & $-1$ & $+1$ \\
$S_6$ & $-1$ & $-1$ \\
$S_7$ & $-1$ & $+i$ \\
$S_8$ & $-1$ & $-i$ \\
\hline
$S_9$ & $+i$ & $+1$ \\
$S_{10}$ & $+i$ & $-1$ \\
$S_{11}$ & $+i$ & $+i$ \\
$S_{12}$ & $+i$ & $-i$ \\
\hline
$S_{13}$ & $-i$ & $+1$ \\
$S_{14}$ & $-i$ & $-1$ \\
$S_{15}$ & $-i$ & $+i$ \\
$S_{16}$ & $-i$ & $-i$ \\
\hline
\end{tabular}
\end{center}

\textbf{Paso 5: Transiciones}

Desde cada estado salen $M = 4$ transiciones (una por cada valor posible de $a_n$).

Por ejemplo, desde estado $S_1 = (+1, +1)$:
\begin{itemize}
    \item Si $a_n = +1$: nuevo estado $(+1, +1) = S_1$
    \item Si $a_n = -1$: nuevo estado $(-1, +1) = S_5$
    \item Si $a_n = +i$: nuevo estado $(+i, +1) = S_9$
    \item Si $a_n = -i$: nuevo estado $(-i, +1) = S_{13}$
\end{itemize}

Valor esperado:
\[
y_{esperado} = h_0 a_n + h_1(+1) + h_2(+1)
\]

\textbf{Paso 6: Complejidad}

\begin{itemize}
    \item \textbf{Estados:} $M^L = 16$
    \item \textbf{Ramas totales por etapa:} $M^L \times M = 16 \times 4 = 64$
    \item \textbf{Operaciones por etapa:} 64 c\'alculos de m\'etrica + 16 comparaciones
\end{itemize}

Comparado con el caso anterior ($L=1$): complejidad se multiplica por 4.

\textbf{Paso 7: Generalizaci\'on}

Para un canal con memoria $L$ y modulaci\'on $M$-aria:
\[
\boxed{\text{N\'umero de estados} = M^L}
\]

\textbf{Ejemplos:}

\begin{center}
\begin{tabular}{|c|c|c|}
\hline
\textbf{Modulaci\'on} & \textbf{Memoria $L$} & \textbf{Estados} \\
\hline
BPSK ($M=2$) & $L=1$ & $2^1 = 2$ \\
BPSK ($M=2$) & $L=2$ & $2^2 = 4$ \\
BPSK ($M=2$) & $L=3$ & $2^3 = 8$ \\
\hline
QPSK ($M=4$) & $L=1$ & $4^1 = 4$ \\
QPSK ($M=4$) & $L=2$ & $4^2 = 16$ \\
QPSK ($M=4$) & $L=3$ & $4^3 = 64$ \\
\hline
16-QAM ($M=16$) & $L=1$ & $16^1 = 16$ \\
16-QAM ($M=16$) & $L=2$ & $16^2 = 256$ \\
\hline
\end{tabular}
\end{center}

\textbf{Implicaci\'on pr\'actica:}

Para canales con ISI larga, el ecualizador ML se vuelve computacionalmente prohibitivo. Por ejemplo:
\begin{itemize}
    \item $L = 5$, QPSK: $4^5 = 1024$ estados
    \item $L = 10$, QPSK: $4^{10} = 1,048,576$ estados
\end{itemize}

\textbf{Soluciones para ISI larga:}

\begin{enumerate}
    \item \textbf{MLSE con estados reducidos (RSSE):} Podar estados poco probables
    \item \textbf{DFE:} Complejidad lineal, pero sub\'optimo
    \item \textbf{OFDM:} Evitar ISI en tiempo usando multi-portadoras
    \item \textbf{Pre-ecualizador ZF/MMSE + ML simplificado:} Reducir $L$ efectiva
\end{enumerate}

\textbf{Conclusi\'on:}

Para el canal $y_n = h_0 a_n + h_1 a_{n-1} + h_2 a_{n-2} + b_n$ con QPSK, el trellis tiene \textbf{16 estados}, lo que representa un aumento significativo en complejidad respecto al caso de memoria $L=1$ (4 estados).

\newpage
\section*{Ejercicio 2 - Sistema OFDM}

\textbf{Contexto del problema:}

Sistema OFDM con las siguientes caracter\'isticas:
\begin{itemize}
    \item Banda total disponible: $B_{total} = 4096$ kHz = $4.096$ MHz
    \item N\'umero total de sub-portadoras: $N = 64$
    \item C\'odigo correcteur de errores: Rendimiento $R_c = 1/2$
    \item Duraci\'on del prefijo c\'iclico: $T_g = T_u/8$ donde $T_u$ es la duraci\'on \'util del s\'imbolo OFDM
    \item Sub-portadoras dedicadas a pilotos: $1/8$ del total
    \item Modulaci\'on por sub-portadora: PSK-8 (8-PSK)
\end{itemize}

\textbf{Relaciones clave:}
\begin{itemize}
    \item Tiempo total de s\'imbolo OFDM: $T_{total} = T_u + T_g = T_u(1 + 1/8) = \frac{9T_u}{8}$
    \item Separaci\'on entre sub-portadoras: $\Delta f = \frac{1}{T_u}$
    \item N\'umero de sub-portadoras de informaci\'on: $N_{info} = N \times (1 - 1/8) = N \times 7/8$
\end{itemize}

\subsection*{Pregunta 1: D\'ebit binaire \'util del sistema}

\textbf{Enunciado:} Calcular el d\'ebit binaire \'util (solo bits de informaci\'on) del sistema.

\vspace{0.3cm}
\textbf{Soluci\'on paso a paso:}

\textbf{Paso 1: Calcular la separaci\'on entre sub-portadoras}

La banda total se reparte entre las $N$ sub-portadoras:
\[
\Delta f = \frac{B_{total}}{N} = \frac{4096 \text{ kHz}}{64} = 64 \text{ kHz}
\]

\textbf{Paso 2: Calcular la duraci\'on \'util $T_u$}

Por definici\'on de OFDM:
\[
T_u = \frac{1}{\Delta f} = \frac{1}{64 \times 10^3} = 15.625 \times 10^{-6} \text{ s} = 15.625 \text{ }\mu\text{s}
\]

\textbf{Paso 3: Calcular el tiempo total de s\'imbolo OFDM}

\[
T_{total} = T_u + T_g = T_u + \frac{T_u}{8} = T_u \left(1 + \frac{1}{8}\right) = \frac{9T_u}{8}
\]

\[
T_{total} = \frac{9 \times 15.625}{8} = 17.578125 \text{ }\mu\text{s}
\]

\textbf{Paso 4: Calcular el n\'umero de sub-portadoras de informaci\'on}

Total de sub-portadoras: $N = 64$

Sub-portadoras para pilotos: $N_{pilot} = \frac{N}{8} = \frac{64}{8} = 8$

Sub-portadoras de informaci\'on: $N_{info} = N - N_{pilot} = 64 - 8 = 56$

\textbf{Paso 5: Calcular bits por s\'imbolo OFDM (antes del c\'odigo)}

Cada sub-portadora transmite un s\'imbolo PSK-8, que transporta:
\[
\log_2(8) = 3 \text{ bits por s\'imbolo}
\]

Bits totales por s\'imbolo OFDM (sub-portadoras de informaci\'on):
\[
B_{OFDM} = N_{info} \times 3 = 56 \times 3 = 168 \text{ bits}
\]

\textbf{Paso 6: Aplicar el rendimiento del c\'odigo corrector}

\textit{\textbf{¿Qu\'e hace el c\'odigo corrector?}} Agrega redundancia para detectar/corregir errores.

Con un c\'odigo de rendimiento $R_c = 1/2$:
\begin{itemize}
    \item Por cada 1 bit de informaci\'on, transmitimos 2 bits (1 bit original + 1 bit de paridad)
    \item Solo la mitad de los bits transmitidos son informaci\'on \'util
    \item La otra mitad es redundancia para protecci\'on
\end{itemize}

\[
B_{info} = B_{OFDM} \times R_c = 168 \times \frac{1}{2} = 84 \text{ bits de informaci\'on por s\'imbolo OFDM}
\]

\textit{\textbf{Trade-off fundamental:}} Sacrificamos 50\% del d\'ebit a cambio de robustez frente a errores.

\textbf{Paso 7: Calcular el d\'ebit binario \'util}

\[
D_b = \frac{B_{info}}{T_{total}} = \frac{84}{17.578125 \times 10^{-6}}
\]

\[
D_b = 4.7808 \times 10^6 \text{ bits/s}
\]

\[
\boxed{D_b = 4.78 \text{ Mbps}}
\]

\textbf{Resumen del c\'alculo:}

\[
D_b = \frac{N_{info} \times \log_2(M) \times R_c}{T_u(1 + T_g/T_u)} = \frac{56 \times 3 \times 0.5}{15.625 \times 10^{-6} \times 1.125} = 4.78 \text{ Mbps}
\]

\textbf{Interpretaci\'on:}

\begin{itemize}
    \item Sin c\'odigo y sin pilotos: $\frac{64 \times 3}{17.578 \times 10^{-6}} = 10.92$ Mbps
    \item Con pilotos (p\'erdida 12.5\%): $10.92 \times 0.875 = 9.56$ Mbps
    \item Con c\'odigo $R=1/2$ (p\'erdida 50\%): $9.56 \times 0.5 = 4.78$ Mbps
\end{itemize}

El c\'odigo correcteur reduce el d\'ebit a cambio de robustez frente a errores.

\subsection*{Pregunta 2: \'Etalement temporel m\'aximo del canal compatible}

\textbf{Enunciado:} Calcular el \'etalement temporel (delay spread) m\'aximo del canal de propagaci\'on compatible con este sistema.

\vspace{0.3cm}
\textbf{Soluci\'on paso a paso:}

\textbf{Paso 1: Condici\'on para evitar ISI entre s\'imbolos OFDM}

\textit{\textbf{¿Qu\'e es el \'etalement temporal (delay spread)?}} Es la diferencia de tiempo entre el primer eco y el \'ultimo eco que llega.

\textit{\textbf{¿Por qu\'e causa problemas?}} Los ecos del s\'imbolo anterior pueden "contaminar" el s\'imbolo actual.

\textit{\textbf{Soluci\'on: El prefijo c\'iclico}} Actuando como "buffer de seguridad":

Condici\'on fundamental:
\[
T_g \geq \tau_{max}
\]

\textit{Interpretaci\'on f\'isica:}
\begin{itemize}
    \item $T_g$ = tiempo del prefijo c\'iclico (copia del final al inicio)
    \item $\tau_{max}$ = retardo m\'aximo del canal (tiempo que tarda el eco m\'as lento)
    \item Si $T_g \geq \tau_{max}$: todos los ecos llegan durante el prefijo, no durante la parte \'util
\end{itemize}

\textit{\textbf{Analog\'{\i}a:}} Es como dejar un margen de tiempo entre trenes para que los pasajeros rezagados no interfieran con el siguiente tren.

\textbf{Paso 2: Calcular $T_g$}

Tenemos:
\[
T_g = \frac{T_u}{8} = \frac{15.625 \text{ }\mu\text{s}}{8} = 1.953125 \text{ }\mu\text{s}
\]

\textbf{Paso 3: Determinar $\tau_{max}$}

Para que el sistema funcione correctamente:
\[
\boxed{\tau_{max} \leq T_g = 1.953125 \text{ }\mu\text{s} \approx 1.95 \text{ }\mu\text{s}}
\]

\textbf{Interpretaci\'on f\'isica:}

En t\'erminos de ancho de banda de coherencia del canal:
\[
B_c = \frac{1}{\tau_{max}} \geq \frac{1}{T_g} = \frac{1}{1.953 \times 10^{-6}} = 512 \text{ kHz}
\]

Como $B_c = 512$ kHz es menor que la banda total (4096 kHz), el canal es \textbf{selectivo en frecuencia}, justificando el uso de OFDM.

\textbf{Ejemplo pr\'actico:}

Si la velocidad de propagaci\'on es $c = 3 \times 10^8$ m/s (ondas electromagn\'eticas en el vac\'io), la diferencia m\'axima de caminos ser\'ia:
\[
\Delta d = c \times \tau_{max} = 3 \times 10^8 \times 1.953 \times 10^{-6} = 586 \text{ metros}
\]

\textbf{Conclusi\'on:}

El sistema puede tolerar canales con \'etalement temporel de hasta $\tau_{max} = 1.95$ μs, correspondiente a diferencias de camino de aproximadamente 586 metros (para ondas de radio) o 2.93 metros (para ondas ac\'usticas submarinas a 1500 m/s).

\subsection*{Pregunta 3: Coh\'erence temporelle minimale del canal compatible}

\textbf{Enunciado:} Calcular la coherencia temporal m\'inima del canal de propagaci\'on compatible con este sistema.

\vspace{0.3cm}
\textbf{Soluci\'on paso a paso:}

\textbf{Paso 1: Condici\'on de coherencia temporal}

\textit{\textbf{¿Qu\'e es el tiempo de coherencia $T_c$?}} Es el tiempo durante el cual el canal permanece "casi constante".

\textit{\textbf{¿Por qu\'e var\'{\i}a el canal?}}
\begin{itemize}
    \item Movimiento del receptor/transmisor (efecto Doppler)
    \item Objetos m\'oviles en el entorno (coches, personas)
    \item Cambios atmosf\'ericos
\end{itemize}

Condici\'on necesaria para OFDM:
\[
T_c \geq T_{total}
\]

\textit{\textbf{¿Por qu\'e esta condici\'on?}} El receptor estima el canal al inicio del s\'imbolo OFDM. Si el canal cambia significativamente durante $T_{total}$, esta estimaci\'on queda obsoleta.

\textit{\textbf{Consecuencia si $T_c < T_{total}$:}} Interferencia entre sub-portadoras (ICI - Inter-Carrier Interference):
\begin{itemize}
    \item Las sub-portadoras pierden su ortogonalidad
    \item La energ\'{\i}a de una sub-portadora "se derrama" a otras
    \item Degradaci\'on severa del rendimiento
\end{itemize}

\textit{\textbf{Analog\'{\i}a:}} Si intentas fotografiar un objeto en movimiento con tiempo de exposici\'on largo, la imagen sale borrosa. Aqu\'{\i} $T_{total}$ es el "tiempo de exposici\'on" del receptor.

\textbf{Paso 2: Calcular $T_{total}$}

Ya calculamos:
\[
T_{total} = T_u + T_g = \frac{9T_u}{8} = 17.578125 \text{ }\mu\text{s}
\]

\textbf{Paso 3: Determinar $T_c$ m\'inimo}

\[
\boxed{T_c \geq T_{total} = 17.58 \text{ }\mu\text{s}}
\]

\textbf{Paso 4: Interpretar en t\'erminos de Doppler}

\textit{\textbf{Relaci\'on fundamental:}} El tiempo de coherencia est\'a relacionado con el m\'aximo desplazamiento Doppler:
\[
T_c \approx \frac{1}{2f_D}
\]

donde $f_D$ es el m\'aximo desplazamiento Doppler.

\textit{\textbf{¿Qu\'e significa?}} Si tengo $T_c = 17.578$ μs, entonces:
\[
f_D \leq \frac{1}{2T_c} = \frac{1}{2 \times 17.578 \times 10^{-6}} \approx 28.44 \text{ kHz}
\]

\textbf{Ejemplo con frecuencia portadora $f_c = 2$ GHz:}

El desplazamiento Doppler se relaciona con la velocidad relativa:
\[
f_D = \frac{v}{c} f_c
\]

Donde $v$ es la velocidad relativa y $c = 3 \times 10^8$ m/s.

Velocidad m\'axima tolerable:
\[
v_{max} = \frac{f_D \cdot c}{f_c} = \frac{28440 \times 3 \times 10^8}{2 \times 10^9} \approx 4266 \text{ m/s} \approx 15360 \text{ km/h}
\]

\textit{\textbf{Interpretaci\'on:}} Este sistema puede tolerar velocidades muy altas (aviones comerciales vuelan a ~900 km/h). En aplicaciones de comunicaciones m\'oviles t\'ipicas (<200 km/h), esta condici\'on se cumple holgadamente.

Donde $v$ es la velocidad relativa y $c = 3 \times 10^8$ m/s.

Velocidad m\'axima tolerable:
\[
v_{max} = \frac{f_D \cdot c}{f_c} = \frac{28440 \times 3 \times 10^8}{2 \times 10^9} \approx 4266 \text{ m/s} \approx 15360 \text{ km/h}
\]

\textit{\textbf{Interpretaci\'on:}} Este sistema puede tolerar velocidades muy altas (aviones comerciales vuelan a ~900 km/h). En aplicaciones de comunicaciones m\'oviles t\'ipicas (<200 km/h), esta condici\'on se cumple holgadamente.

\textbf{Paso 5: Verificaci\'on pr\'actica}

\textit{\textbf{Ejemplo realista:}} Usuario en un veh\'iculo a 100 km/h

Velocidad: $v = 100$ km/h = 27.78 m/s

Desplazamiento Doppler:
\[
f_D = \frac{v}{c} \times f_c = \frac{27.78}{3 \times 10^8} \times 2 \times 10^9 = 185 \text{ Hz}
\]

Tiempo de coherencia del canal:
\[
T_c \approx \frac{1}{2 \times 185} = 2.7 \text{ ms} = 2700 \text{ }\mu\text{s}
\]

\textit{\textbf{Comparaci\'on:}} $2700$ μs (canal a 100 km/h) $\gg 17.58$ μs (requerimiento del sistema)

\textcolor{green}{\textbf{Conclusi\'on:}} El sistema es \textbf{perfectamente compatible} con este escenario m\'ovil. Hay un margen de seguridad de $2700/17.58 \approx 154$ veces.

\textbf{Paso 6: Conclusi\'on general}

El sistema requiere un tiempo de coherencia m\'inimo de $\boxed{T_c \geq 17.58 \text{ }\mu\text{s}}$. 

\textit{Esto es compatible con:}
\begin{itemize}
    \item \textbf{Canales fijos:} interior de edificios, enlaces punto a punto ($T_c$ casi infinito)
    \item \textbf{Movilidad moderada:} hasta velocidades de ~15,000 km/h (mucho m\'as r\'apido que cualquier veh\'iculo terrestre)
    \item \textbf{Canales variables lentos:} cambios atmosf\'ericos, movimiento de objetos en el entorno
\end{itemize}

\textit{NO es compatible con:}
\begin{itemize}
    \item Canales de alta movilidad relativa a frecuencias muy altas
    \item Escenarios con scattering muy r\'apido
\end{itemize}

\newpage
\section*{Ejercicio 3 - C\'odigo a repetici\'on}

\textit{\textbf{Contexto del ejercicio:}} Estudiaremos el c\'odigo corrector de errores m\'as simple: el \textbf{c\'odigo de repetici\'on (3,1)}. Este c\'odigo transmite cada bit 3 veces y usa votaci\'on por mayor\'ia para decidir. Aunque simple, ilustra conceptos fundamentales de teor\'ia de c\'odigos.

\textit{\textbf{Par\'ametros del c\'odigo:}}
\begin{itemize}
    \item $n = 3$: longitud de la palabra c\'odigo (bits transmitidos)
    \item $k = 1$: longitud del mensaje (bits de informaci\'on)  
    \item Rendimiento: $R = k/n = 1/3 \approx 33.3\%$
    \item Distancia m\'inima: $d_{min} = 3$ (palabras $(0,0,0)$ y $(1,1,1)$ difieren en 3 posiciones)
    \item Capacidad de correcci\'on: $t = \lfloor (d_{min}-1)/2 \rfloor = 1$ error
\end{itemize}

\textbf{Contexto del problema:}

\begin{itemize}
    \item Probabilidad de error binaria del canal sin c\'odigo: $p_e$
    \item C\'odigo utilizado: C\'odigo de repetici\'on de rendimiento $R = 1/3$
    \item Esquema: Cada bit se repite 3 veces antes de transmitir
    \item Regla de decodificaci\'on: Decisi\'on por mayor\'ia
\end{itemize}

\subsection*{Pregunta 1: Probabilidad de error despu\'es de la decodificaci\'on}

\textbf{Enunciado:} Calcular la probabilidad de error binaria despu\'es de la decodificaci\'on.

\vspace{0.3cm}
\textbf{Soluci\'on paso a paso:}

\textbf{Paso 1: Funcionamiento del c\'odigo de repetici\'on}

\textit{\textbf{Idea simple pero poderosa:}} "Si algo es importante, d\'ilo tres veces."

\textbf{En el emisor:}
\begin{itemize}
    \item Bit de informaci\'on original: $b \in \{0, 1\}$
    \item Palabra c\'odigo transmitida: $(b, b, b)$ - repetimos 3 veces
    \item Ejemplo: $b=0 \to$ transmitimos $(0, 0, 0)$
    \item Ejemplo: $b=1 \to$ transmitimos $(1, 1, 1)$
\end{itemize}

\textit{Costo:} Usamos 3 bits de canal para 1 bit de informaci\'on $\to$ Rendimiento $R = 1/3$

\textbf{En el canal (ruidoso):}
\begin{itemize}
    \item Cada bit tiene probabilidad $p_e$ de invertirse (error)
    \item Los errores son independientes entre s\'i
    \item Palabra recibida: $(r_1, r_2, r_3)$ con posibles errores
    \item Ejemplo: transmitimos $(0,0,0)$, puede llegar $(0,1,0)$ si hay error en bit 2
\end{itemize}

\textbf{En el decodificador (receptor):}
\begin{itemize}
    \item \textbf{Regla democr\'atica - Votaci\'on por mayor\'ia:}
    \[
    \hat{b} = \begin{cases}
    1 & \text{si al menos 2 de } \{r_1, r_2, r_3\} \text{ son 1} \\
    0 & \text{en caso contrario}
    \end{cases}
    \]
    \item Ejemplo 1: recibimos $(0,0,1)$ $\to$ mayor\'ia de 0's $\to$ decidimos $\hat{b}=0$ ✓
    \item Ejemplo 2: recibimos $(0,1,1)$ $\to$ mayor\'ia de 1's $\to$ decidimos $\hat{b}=1$
\end{itemize}

\textit{\textbf{Principio:}} Podemos corregir hasta 1 error, porque la mayor\'ia sigue siendo correcta. Pero si hay 2 o 3 errores, perdemos.

\textbf{Paso 2: An\'alisis de error}

Supongamos que se transmite $b = 0$, es decir, la palabra c\'odigo es $(0, 0, 0)$.

\textbf{Error ocurre si:} Al menos 2 de los 3 bits recibidos son err\'oneos (se convierten en 1).

\textbf{Casos de error:}
\begin{enumerate}
    \item Exactamente 2 bits err\'oneos: $(1, 1, 0)$, $(1, 0, 1)$, $(0, 1, 1)$
    \item Los 3 bits err\'oneos: $(1, 1, 1)$
\end{enumerate}

\textbf{Paso 3: C\'alculo de probabilidades}

\textbf{Probabilidad de exactamente 2 errores:}

N\'umero de combinaciones: $\binom{3}{2} = 3$

Probabilidad:
\[
P(2 \text{ errores}) = \binom{3}{2} p_e^2 (1-p_e) = 3p_e^2(1-p_e)
\]

\textbf{Probabilidad de 3 errores:}

\[
P(3 \text{ errores}) = p_e^3
\]

\textbf{Paso 4: Probabilidad total de error despu\'es de decodificaci\'on}

\[
P_e^{(decod)} = P(2 \text{ errores}) + P(3 \text{ errores})
\]

\[
P_e^{(decod)} = 3p_e^2(1-p_e) + p_e^3
\]

\[
P_e^{(decod)} = 3p_e^2 - 3p_e^3 + p_e^3
\]

\[
\boxed{P_e^{(decod)} = 3p_e^2 - 2p_e^3}
\]

\textbf{Forma alternativa:}

\[
P_e^{(decod)} = p_e^2(3 - 2p_e)
\]

\textbf{Paso 5: An\'alisis de ganancia del c\'odigo}

\textbf{Comparaci\'on con canal sin c\'odigo:}

\begin{itemize}
    \item Sin c\'odigo: $P_e = p_e$
    \item Con c\'odigo: $P_e^{(decod)} = 3p_e^2 - 2p_e^3$
\end{itemize}

\textbf{El c\'odigo es beneficioso si:}
\[
3p_e^2 - 2p_e^3 < p_e
\]

Dividiendo por $p_e$ (asumiendo $p_e > 0$):
\[
3p_e - 2p_e^2 < 1
\]

\[
2p_e^2 - 3p_e + 1 > 0
\]

Factorizando:
\[
(2p_e - 1)(p_e - 1) > 0
\]

Soluciones: $p_e < 1/2$ o $p_e > 1$ (imposible)

\textbf{Conclusi\'on:} El c\'odigo de repetici\'on es efectivo si $\boxed{p_e < 0.5}$

\textbf{Paso 6: Ejemplos num\'ericos de ganancia}

\textit{\textbf{Escenario 1:}} Canal con $p_e = 0.1$ (BER del 10\%)

Sin c\'odigo:
\[
P_e = 0.1 = 10\%
\]

Con c\'odigo de repetici\'on:
\[
P_e^{(decod)} = 3(0.1)^2 - 2(0.1)^3 = 0.03 - 0.002 = 0.028 = 2.8\%
\]

\textit{\textbf{Ganancia:}} $10\% \to 2.8\%$ - Reducci\'on del $72\%$ en la tasa de error.

\vspace{0.2cm}

\textit{\textbf{Escenario 2:}} Canal con $p_e = 0.01$ (BER del 1\%)

Sin c\'odigo:
\[
P_e = 0.01 = 1\%
\]

Con c\'odigo de repetici\'on:
\[
P_e^{(decod)} = 3(0.01)^2 - 2(0.01)^3 = 0.0003 - 0.000002 \approx 0.0003 = 0.03\%
\]

\textit{\textbf{Ganancia espectacular:}} $1\% \to 0.03\%$ - Reducci\'on del $97\%$ en la tasa de error.

\vspace{0.2cm}

\textit{\textbf{Observaci\'on clave:}} Cuanto mejor sea el canal (menor $p_e$), mayor es la ganancia relativa del c\'odigo.

\vspace{0.2cm}

\textit{\textbf{Escenario 3:}} Canal malo con $p_e = 0.3$ (BER del 30\%)

Sin c\'odigo:
\[
P_e = 0.3 = 30\%
\]

Con c\'odigo de repetici\'on:
\[
P_e^{(decod)} = 3(0.3)^2 - 2(0.3)^3 = 0.27 - 0.054 = 0.216 = 21.6\%
\]

\textit{Ganancia modesta:} $30\% \to 21.6\%$ - Reducci\'on del $28\%$ solamente.

\vspace{0.2cm}

\textit{\textbf{Conclusi\'on general:}} El c\'odigo de repetici\'on funciona mejor en canales con BER bajo. Para canales muy ruidosos ($p_e$ cercano a 0.5), la ganancia es marginal.

\textbf{Ejemplos num\'ericos:}

\begin{center}
\begin{tabular}{|c|c|c|c|}
\hline
\textbf{$p_e$} & \textbf{$P_e^{(decod)}$} & \textbf{Ganancia} & \textbf{Factor de mejora} \\
\hline
0.1 & $3(0.01) - 2(0.001) = 0.028$ & Mejora 3.6$\times$ & $p_e/P_e = 3.57$ \\
0.01 & $3(10^{-4}) - 2(10^{-6}) \approx 2.98 \times 10^{-4}$ & Mejora 33.6$\times$ & $p_e/P_e = 33.56$ \\
0.001 & $3(10^{-6}) - 2(10^{-9}) \approx 3 \times 10^{-6}$ & Mejora 333$\times$ & $p_e/P_e = 333$ \\
0.5 & $3(0.25) - 2(0.125) = 0.5$ & Sin mejora & Factor = 1 \\
\hline
\end{tabular}
\end{center}

\textbf{Observaci\'on:} Para $p_e$ peque\~nos, $P_e^{(decod)} \approx 3p_e^2 \ll p_e$.

\subsection*{Pregunta 2: Matriz de paridad, lista de s\'indromes y motivos de error}

\textbf{Enunciado:} Dar la matriz de paridad de este c\'odigo, la lista de s\'indromes y los motivos de error asociados a cada valor de s\'indrome.

\vspace{0.3cm}
\textbf{Soluci\'on paso a paso:}

\textbf{Paso 1: Par\'ametros del c\'odigo}

\textit{\textbf{¿Qu\'e es un c\'odigo lineal?}} Es un c\'odigo donde las palabras v\'alidas forman un subespacio vectorial.

\begin{itemize}
    \item C\'odigo de repetici\'on $(n, k) = (3, 1)$
    \item $n = 3$: Longitud de palabra c\'odigo (bits transmitidos)
    \item $k = 1$: N\'umero de bits de informaci\'on (mensaje)
    \item $r = n - k = 2$: Bits de paridad (redundancia)
\end{itemize}

\textbf{Paso 2: Matriz generatriz $G$}

\textit{\textbf{¿Qu\'e hace la matriz $G$?}} Transforma el mensaje en la palabra c\'odigo.

El c\'odigo de repetici\'on tiene matriz generatriz:
\[
G = \begin{bmatrix} 1 & 1 & 1 \end{bmatrix}
\]

\textit{Operaci\'on de codificaci\'on:}
\[
\mathbf{c} = b \cdot G = b \times [1, 1, 1] = [b, b, b]
\]

\textit{Ejemplos:}
\begin{itemize}
    \item Si $b = 0$: $\mathbf{c} = [0, 0, 0]$
    \item Si $b = 1$: $\mathbf{c} = [1, 1, 1]$
\end{itemize}

\textit{\textbf{Observaci\'on:}} Solo hay 2 palabras c\'odigo v\'alidas (de las $2^3 = 8$ posibles).

\textbf{Paso 3: Matriz de paridad $H$}

\textit{\textbf{¿Qu\'e hace la matriz $H$?}} Detecta si una palabra recibida es v\'alida (sin errores) o no.

\textit{\textbf{Propiedad fundamental:}} Para cualquier palabra c\'odigo v\'alida $\mathbf{c}$:
\[
H \cdot \mathbf{c}^T = \mathbf{0}
\]

La matriz de paridad debe satisfacer: $G \cdot H^T = \mathbf{0}$ (ortogonalidad).

Para el c\'odigo de repetici\'on $(3,1)$:
\[
\boxed{H = \begin{bmatrix}
1 & 1 & 0 \\
1 & 0 & 1
\end{bmatrix}}
\]

\textit{Dimensiones:} $H$ es $(n-k) \times n = 2 \times 3$

\textbf{Verificaci\'on de ortogonalidad:}
\[
G \cdot H^T = \begin{bmatrix} 1 & 1 & 1 \end{bmatrix} \cdot \begin{bmatrix}
1 & 1 \\
1 & 0 \\
0 & 1
\end{bmatrix} = \begin{bmatrix} 1+1+0 & 1+0+1 \end{bmatrix}
\]

En aritm\'etica binaria (m\'odulo 2):
\[
= \begin{bmatrix} 0 & 0 \end{bmatrix} \quad \checkmark
\]

\textit{\textbf{¿Por qu\'e funciona?}} Las filas de $H$ definen "ecuaciones de paridad" que toda palabra v\'alida debe satisfacer:
\begin{itemize}
    \item Fila 1: $c_1 \oplus c_2 = 0$ (los primeros dos bits deben ser iguales)
    \item Fila 2: $c_1 \oplus c_3 = 0$ (el primero y tercero deben ser iguales)
\end{itemize}

Si ambas se cumplen, entonces $c_1 = c_2 = c_3$ (repetici\'on).

\textbf{Paso 4: C\'alculo del s\'indrome}

\textit{\textbf{¿Qu\'e es el s\'indrome?}} Es un vector que nos dice si hay error y, en caso afirmativo, d\'onde est\'a.

Para una palabra recibida $\mathbf{r} = [r_1, r_2, r_3]$:
\[
\mathbf{s} = H \cdot \mathbf{r}^T = \begin{bmatrix}
1 & 1 & 0 \\
1 & 0 & 1
\end{bmatrix} \cdot \begin{bmatrix}
r_1 \\
r_2 \\
r_3
\end{bmatrix} = \begin{bmatrix}
r_1 \oplus r_2 \\
r_1 \oplus r_3
\end{bmatrix}
\]

donde $\oplus$ denota suma m\'odulo 2 (XOR).

\textit{\textbf{Interpretaci\'on de cada componente:}}
\begin{itemize}
    \item $s_1 = r_1 \oplus r_2$: Compara bits 1 y 2. Si son diferentes, $s_1 = 1$
    \item $s_2 = r_1 \oplus r_3$: Compara bits 1 y 3. Si son diferentes, $s_2 = 1$
\end{itemize}

\textbf{Paso 5: Lista de s\'indromes y motivos de error}

\textit{\textbf{Principio del s\'indrome:}} Cada patr\'on de error de un solo bit produce un s\'indrome \\'unico.

\begin{center}
\begin{tabular}{|c|c|c|c|c|}
\hline
\textbf{S\'indrome $\mathbf{s}$} & \textbf{Motivo $\mathbf{e}$} & \textbf{Descripci\'on} & \textbf{Explicaci\'on} & \textbf{Acci\'on} \\
\hline
$[0, 0]$ & $[0, 0, 0]$ & Sin error & $r_1 = r_2 = r_3$ & No corregir \\
\hline
$[1, 1]$ & $[1, 0, 0]$ & Error en bit 1 & $r_1 \neq r_2$ y $r_1 \neq r_3$ & Invertir $r_1$ \\
\hline
$[1, 0]$ & $[0, 1, 0]$ & Error en bit 2 & $r_1 \neq r_2$ pero $r_1 = r_3$ & Invertir $r_2$ \\
\hline
$[0, 1]$ & $[0, 0, 1]$ & Error en bit 3 & $r_1 = r_2$ pero $r_1 \neq r_3$ & Invertir $r_3$ \\
\hline
\end{tabular}
\end{center}

\textbf{Paso 6: Verificaci\'on con ejemplos concretos}

\textit{\textbf{Ejemplo 1:}} Transmitimos $[0,0,0]$ pero llega $[0,1,0]$ (error en bit 2)

C\'alculo del s\'indrome:
\[
\mathbf{s} = \begin{bmatrix}
r_1 \oplus r_2 \\
r_1 \oplus r_3
\end{bmatrix} = \begin{bmatrix}
0 \oplus 1 \\
0 \oplus 0
\end{bmatrix} = \begin{bmatrix}
1 \\
0
\end{bmatrix}
\]

\textit{S\'indrome $[1,0]$ indica error en bit 2.} Acci\'on: invertir $r_2$.

Palabra corregida: $[0,0,0]$ ✓

\vspace{0.3cm}

\textit{\textbf{Ejemplo 2:}} Transmitimos $[1,1,1]$ pero llega $[1,0,1]$ (error en bit 2)

C\'alculo del s\'indrome:
\[
\mathbf{s} = \begin{bmatrix}
1 \oplus 0 \\
1 \oplus 1
\end{bmatrix} = \begin{bmatrix}
1 \\
0
\end{bmatrix}
\]

\textit{S\'indrome $[1,0]$ indica error en bit 2.} Acci\'on: invertir $r_2$.

Palabra corregida: $[1,1,1]$ ✓

\vspace{0.3cm}

\textit{\textbf{Ejemplo 3:}} Transmitimos $[0,0,0]$ pero llega $[1,0,0]$ (error en bit 1)

C\'alculo del s\'indrome:
\[
\mathbf{s} = \begin{bmatrix}
1 \oplus 0 \\
1 \oplus 0
\end{bmatrix} = \begin{bmatrix}
1 \\
1
\end{bmatrix}
\]

\textit{S\'indrome $[1,1]$ indica error en bit 1.} Acci\'on: invertir $r_1$.

Palabra corregida: $[0,0,0]$ ✓

\textbf{Paso 7: Tabla completa de decodificaci\'on}

\begin{center}
\small
\begin{tabular}{|c|c|c|c|c|}
\hline
\textbf{Recibido $\mathbf{r}$} & \textbf{S\'indrome $\mathbf{s}$} & \textbf{Error $\mathbf{e}$} & \textbf{Corregido $\mathbf{c}$} & \textbf{Bit decodificado} \\
\hline
$[0, 0, 0]$ & $[0, 0]$ & $[0, 0, 0]$ & $[0, 0, 0]$ & 0 \\
$[0, 0, 1]$ & $[0, 1]$ & $[0, 0, 1]$ & $[0, 0, 0]$ & 0 \\
$[0, 1, 0]$ & $[1, 0]$ & $[0, 1, 0]$ & $[0, 0, 0]$ & 0 \\
$[0, 1, 1]$ & $[1, 1]$ & $[1, 0, 0]$ & $[1, 1, 1]$ & 1 (mayor\'ia) \\
\hline
$[1, 0, 0]$ & $[1, 1]$ & $[1, 0, 0]$ & $[0, 0, 0]$ & 0 \\
$[1, 0, 1]$ & $[1, 0]$ & $[0, 1, 0]$ & $[1, 1, 1]$ & 1 (mayor\'ia) \\
$[1, 1, 0]$ & $[0, 1]$ & $[0, 0, 1]$ & $[1, 1, 1]$ & 1 (mayor\'ia) \\
$[1, 1, 1]$ & $[0, 0]$ & $[0, 0, 0]$ & $[1, 1, 1]$ & 1 \\
\hline
\end{tabular}
\end{center}

\textbf{Paso 8: Observaci\'on importante sobre limitaciones}

\textit{\textbf{Capacidad del c\'odigo:}} El c\'odigo de repetici\'on $(3,1)$ puede corregir **todos los errores simples** (1 bit err\'oneo).

\textit{\textbf{Problema con errores m\'ultiples:}}

\textit{Caso problem\'atico:} Transmitimos $[0,0,0]$ pero llega $[1,1,0]$ (2 errores en bits 1 y 2)

C\'alculo del s\'indrome:
\[
\mathbf{s} = \begin{bmatrix}
1 \oplus 1 \\
1 \oplus 0
\end{bmatrix} = \begin{bmatrix}
0 \\
1
\end{bmatrix}
\]

El decodificador interpreta $[0,1]$ como error en bit 3 (falso). Correcci\'on: invierte $r_3$.

Resultado: $[1,1,1]$ \textcolor{red}{✗ (INCORRECTO - debimos obtener $[0,0,0]$)}

\vspace{0.2cm}

\textit{\textbf{Conclusi\'on cr\'itica:}}
\begin{itemize}
    \item Errores simples: \textcolor{green}{siempre se corrigen correctamente}
    \item Errores dobles: \textcolor{red}{se detectan pero se corrigen incorrectamente}
    \item Errores triples: \textcolor{red}{pasan desapercibidos (s\'indrome $[0,0]$)}
\end{itemize}

\textit{Principio general:} Un c\'odigo con distancia m\'inima $d_{min} = 3$ puede corregir hasta $t = 1$ error o detectar hasta 2 errores (pero no ambos simult\'aneamente).

Esta es la raz\'on por la cual la probabilidad de error despu\'es de decodificaci\'on es $3p_e^2 - 2p_e^3$ (errores dobles y triples causan error de decodificaci\'on).

\textbf{Conclusi\'on:}

El c\'odigo de repetici\'on $(3,1)$ tiene:
\begin{itemize}
    \item Matriz de paridad $H = \begin{bmatrix} 1 & 1 & 0 \\ 1 & 0 & 1 \end{bmatrix}$
    \item 4 s\'indromes posibles: $[0,0], [1,0], [0,1], [1,1]$
    \item Capacidad de correcci\'on: 1 error
    \item Capacidad de detecci\'on: 2 errores (pero no los corrige)
    \item Distancia m\'inima: $d_{min} = 3$
\end{itemize}

\end{document}
